\documentclass[11pt,a4paper]{article}
\usepackage{oi-statement}
\OISetup{CSP-J 模拟赛}{第七套}{入门级}
\OISetCompileOptions{-O2 -std=c++14}
\begin{document}
\OICoverTitle
\OICoverTable{4}{
题目名称 & 回文串 & 交错四元组 & 文件判定 & 树与叶子\\\hline

题目类型 & 传统型 & 传统型 & 传统型 & 传统型\\\hline

目录 & {\ttfamily palindrome} & {\ttfamily quad} & {\ttfamily files} & {\ttfamily leafcut}\\\hline

可执行文件名 & {\ttfamily palindrome} & {\ttfamily quad} & {\ttfamily files} & {\ttfamily leafcut}\\\hline

输入文件名 & {\ttfamily palindrome.in} & {\ttfamily quad.in} & {\ttfamily files.in} & {\ttfamily leafcut.in}\\\hline

输出文件名 & {\ttfamily palindrome.out} & {\ttfamily quad.out} & {\ttfamily files.out} & {\ttfamily leafcut.out}\\\hline

每个测试点时限 & 1.0 秒 & 1.0 秒 & 1.0 秒 & 1.0 秒\\\hline

内存限制 & 512 MiB & 512 MiB & 512 MiB & 512 MiB\\\hline

测试点数目 & 10 & 10 & 10 & 10\\\hline

测试点是否等分 & 是，每点 10 分 & 是，每点 10 分 & 是，每点 10 分 & 是，每点 10 分\\\hline
}
\OICompilation{4}{
对于 C++ 语言 & {\ttfamily palindrome.cpp} & {\ttfamily quad.cpp} & {\ttfamily files.cpp} & {\ttfamily leafcut.cpp}\\\hline
}
\OICoverNotes
\begin{enumerate}
\item 本套为 CSP-J 训练模拟赛，共 4 题，每题 100 分，总分 400 分。
\item 各题均采用标准输入输出，提交对应的同名 C++ 源程序。
\item 每题有 10 个测试点，每点 10 分，按测试点独立计分。
\item 答案采用普通文本比较，忽略行末空格和文件末尾多余换行。
\item 题面样例与下发大样例不计入正式测试点。
\item 程序可使用的栈空间与题目的内存限制一致。
\item 样例框左侧行号不属于输入输出。
\item 使用 C++14，编译选项为 \texttt{-O2 -std=c++14}。
\end{enumerate}

\OIProblem{回文串}{palindrome}

\OISection{题目描述}

给定一个字符串 $s$。你可以任意重排其中的字符，求重排后回文子串数量的最大值。

子串必须连续且非空。一个字符串正着读与反着读相同，就称为回文串，单个字符也是回文串。

按子串的位置区间计数：即使内容相同，只要起点或终点不同，就视为不同子串。例如 \texttt{\detokenize{aa}} 有两个单字符回文子串和一个长度为 $2$ 的回文子串。

\OISection{输入格式}

输入一行字符串 $s$。

\OISection{输出格式}

输出一个整数，表示最多的回文子串数量。

\begin{OISample}{1}{输入}
abca
\end{OISample}
\begin{OISample}{1}{输出}
5
\end{OISample}
\OISection{样例 1 解释}

可以排列成 \texttt{\detokenize{aabc}}，回文子串为两个 \texttt{\detokenize{a}}、一个 \texttt{\detokenize{b}}、一个 \texttt{\detokenize{c}} 和一个 \texttt{\detokenize{aa}}，共 $5$ 个。

\begin{OISample}{2}{输入}
aa
\end{OISample}
\begin{OISample}{2}{输出}
3
\end{OISample}
\OISection{样例 2 解释}

两个 \texttt{\detokenize{a}} 和一个 \texttt{\detokenize{aa}}，共 $3$ 个。

\OIAdditionalSample{3}{palindrome/palindrome}
\OIAdditionalSample{4}{palindrome/palindrome}
\Needspace{6\baselineskip}
\OISection{数据范围}

$1\le |s|\le10^6$；字符串只包含小写英文字母。

每个测试点 10 分，共 10 个测试点，总分 100 分。

\begin{OIDataTable}{|c|c|c|}
测试点 & 本组分值 & 额外限制\\\hline
1～3 & 30 & $1\le \lvert s\rvert\le3$\\\hline
4～10 & 70 & 无额外限制\\\hline
\end{OIDataTable}

\OIProblem{交错四元组}{quad}

\OISection{题目描述}

给定一个长度为 $n$ 的整数数组 $a$，从左到右选择四个位置 $i<j<k<l$，要求

\[
a_i<a_k<a_j<a_l.
\]

求满足条件的下标四元组数量。四个不等号中的比较均为严格比较。数组可以包含相等元素；不同的下标选择分别计数。

\OISection{输入格式}

第一行一个整数 $n$。第二行 $n$ 个整数 $a_1,\ldots,a_n$。

\OISection{输出格式}

输出一个整数，表示合法选择的数量。

\begin{OISample}{1}{输入}
4
1 3 2 4
\end{OISample}
\begin{OISample}{1}{输出}
1
\end{OISample}
\begin{OISample}{2}{输入}
5
1 2 3 4 5
\end{OISample}
\begin{OISample}{2}{输出}
0
\end{OISample}
\begin{OISample}{3}{输入}
5
1 4 2 3 5
\end{OISample}
\begin{OISample}{3}{输出}
2
\end{OISample}
\OISection{样例 3 解释}

合法下标组为 $(1,2,3,5)$ 和 $(1,2,4,5)$，答案为 $2$。

\OIAdditionalSample{4}{quad/quad}
\OIAdditionalSample{5}{quad/quad}
\Needspace{7\baselineskip}
\OISection{数据范围}

$4\le n\le2000$，$-10^9\le a_i\le10^9$。前两档累计分别为 20\%、50\%。

每个测试点 10 分，共 10 个测试点，总分 100 分。

\begin{OIDataTable}{|c|c|c|}
测试点 & 本组分值 & 额外限制\\\hline
1～2 & 20 & $n\le20$\\\hline
3～5 & 30 & $n\le800$\\\hline
6～10 & 50 & 无额外限制\\\hline
\end{OIDataTable}

\OIProblem{文件判定}{files}

\OISection{题目描述}

一次练习赛有四道题。选手的提交目录必须满足以下规则，请判断给出的全部文件关系是否合法。

\begin{enumerate}
\item 只有一个根目录。根目录名称由一段非空英文字母和一段非空数字依次连接而成，例如 \texttt{\detokenize{BJ01}}；\texttt{\detokenize{123}}、\texttt{\detokenize{ABC}}、\texttt{\detokenize{A2B}} 都不合法。
\item 根目录中恰好有四个题目文件夹，每个题目名称是非空的纯英文字母串。
\item 每个题目文件夹恰好包含一个条目，名称必须为“题目名称 + \texttt{\detokenize{.cpp}}”，并且该条目不能再有子条目。
\item 大小写敏感。不能出现多余文件或多余目录。只判断名称和包含关系，不区分文件与文件夹的真实类型。
\end{enumerate}

例如，合法目录可以是：

\begin{quote}\ttfamily
BJ01
  apple / apple.cpp
  Banana / Banana.cpp
  candy / candy.cpp
  task / task.cpp
\end{quote}

输入中的每条关系 $A\ B$ 表示 $B$ 是 $A$ 的直接子条目。关系不一定按目录层次顺序给出。

\OISection{输入格式}

第一行正整数 $t$，表示数据组数。

每组数据先给出正整数 $n$，表示关系条数。接下来 $n$ 行，每行两个字符串 $A,B$。

同一名称表示同一个条目，保证同名条目不会出现在目录的不同位置；给出的关系不重复，并且 $A\ne B$。

\OISection{输出格式}

每组输出一行 \texttt{\detokenize{yes}} 或 \texttt{\detokenize{no}}。

\begin{OISample}{1}{输入}
5
2
BJ01 task
task task.cpp
8
task task.cpp
bJ01 task
bJ01 apple
apple apple.cpp
bJ01 Banana
Banana Banana.cpp
bJ01 candy
candy candy.cpp
8
123 task
task task.cpp
123 plane
plane plane.cpp
123 banana
banana banana.cpp
123 candy
candy candy.cpp
8
BJ01 task
task task.cpp
BJ01 apple
apple banana.cpp
BJ01 banana
banana apple.cpp
BJ01 candy
candy candy.cpp
8
BJ01 task
task task.cpp
BJ01 apple
apple apple.cpp
BJ01 banana
banana banana.cpp
BJ01 candy
candy candy.cpp.cpp
\end{OISample}
\begin{OISample}{1}{输出}
no
yes
no
no
no
\end{OISample}
\OISection{样例 1 解释}

五组依次为：题目数量不足；合法；根名称没有字母；文件与题目名称不对应；扩展名重复。

\OIAdditionalSample{2}{files/files}
\OIAdditionalSample{3}{files/files}
\Needspace{11\baselineskip}
\OISection{数据范围}

$1\le t\le100$，每组 $1\le n\le10$。每个名称长度为 $1\sim20$，只可能包含英文字母、数字以及 \texttt{\detokenize{- . ( )}}。各行独立计分。

性质 1：根目录唯一且其名称符合规则。

性质 2：所有条目的名称均符合对应要求，只可能在根目录个数或第二级目录个数上出错。

性质 3：目录的最大层数恰好为三级。

每个测试点 10 分，共 10 个测试点，总分 100 分。

\begin{OIDataTable}{|c|c|c|}
测试点 & 本组分值 & 额外限制\\\hline
1～2 & 20 & 同时满足性质 1、2、3\\\hline
3～4 & 20 & 同时满足性质 1、3\\\hline
5～6 & 20 & 满足性质 2\\\hline
7～10 & 40 & 无额外性质\\\hline
\end{OIDataTable}

\OIProblem{树与叶子}{leafcut}

\OISection{题目描述}

给定一棵 $n$ 个点的无根树。一个点的度是与它直接相连的点数。本题将度不超过 $1$ 的点定义为叶子，因此孤立的单个点也是叶子。

删除树中的一条边后，树分成两个连通部分。分别按照删除后的度数，统计两部分的叶子数量 $L_1,L_2$。

求 $|L_1-L_2|$ 的最小值，以及有多少条边在删除后能达到这个最小值。

\OISection{输入格式}

第一行正整数 $T$。每组先输入整数 $n$，接下来 $n-1$ 行，每行两个整数 $u,v$，表示一条无向边。

\OISection{输出格式}

每组输出两个非负整数：最小叶子数量差值，以及达到该差值的断边方案数。

\begin{OISample}{1}{输入}
2
4
1 2
2 3
2 4
4
2 1
1 3
3 4
\end{OISample}
\begin{OISample}{1}{输出}
1 3
0 1
\end{OISample}
\OISection{样例 1 解释}

第一组无论删除哪条边，都是一个孤点和一棵三点小树，叶子数分别为 $1,2$，所以输出 \texttt{\detokenize{1 3}}。第二组是一条四点链，只有删去中间边能让两边各有两个叶子。

\OIAdditionalSample{2}{leafcut/leafcut}
\OIAdditionalSample{3}{leafcut/leafcut}
\Needspace{9\baselineskip}
\OISection{数据范围}

$1\le T\le10$，每组 $2\le n\le10^5$，$1\le u,v\le n$。保证输入为一棵树。前两档累计为 10\%、20\%；链和方案数为 1 的两档各另占 10\%。

每个测试点 10 分，共 10 个测试点，总分 100 分。

\begin{OIDataTable}{|c|c|c|}
测试点 & 本组分值 & 额外限制\\\hline
1 & 10 & $n\le100$，$T=1$\\\hline
2 & 10 & $n\le1000$\\\hline
3 & 10 & 树是一条链\\\hline
4 & 10 & 达到最小差值的断边方案数为 1\\\hline
5～10 & 60 & 无额外限制\\\hline
\end{OIDataTable}

\end{document}
